\chapter{周期：为什么好日子和坏日子会轮流来}
\chaptermeta{18}

\begin{ledetext}
你算不出下一个拐点，这件事没有例外。但你能判断自己现在站在哪一段，而这才是有用的那个问题。
\end{ledetext}

\section{先把钟表扔了}

"下一个拐点什么时候到"是这本书里我可以最干脆回答的问题：不知道，而且没有人知道。

"周期"这个词本身就在骗人。它让人想起钟摆和正弦曲线，想起"上一轮走了 8 年，所以这一轮也是 8 年"。真实的周期长度从三年到三十年不等，方差大得离谱，而且只能事后量出来。

能做的只有一件事：判断自己大概站在什么位置。不是问"什么时候转向"，是问"现在算早、算中、还是算晚"。

这两个问题的难度差着一个数量级。

像在山里判断天气。你说不出下午三点几分下雨，但你能看云、看风向、看气压计，判断出"今天大概率要下"，于是带伞、不上山脊、下午四点前撤。做不到预测，做得到不被淋死。

这个比喻还能再推一步：山里真正的失误从来不是"没带伞"，是明明看见云压下来了，因为已经爬了三个小时不甘心，硬要往上走。周期里也一样。绝大部分亏损不来自看错，来自不肯认。

\section{猪}

河南南阳有个养猪场，场主姓宋，2019 年的时候存栏三千头。

那年非洲猪瘟叠加产能出清，猪价一路冲到 40 块一公斤。2020 年出栏一头猪，毛利两千多。宋老板做了一个所有人都会做的决定：扩栏，从 3000 头加到 8000 头。

问题在于从决定到出肉要多久。母猪配种到产仔 114 天，仔猪育肥到出栏 6 个月上下。如果是从买后备母猪开始算，完整一轮接近 14 到 18 个月。

他 2020 年做的决定，2021 年下半年才变成猪肉。

而全国所有看到同一个价格的养殖户，在同一段时间做了同一个决定。

2021 年 10 月，生猪价格跌到 10 块出头一公斤，全行业亏损。宋老板开始淘汰母猪。全国都在淘汰母猪。淘汰的效果同样要等一年多才显现，于是 2022 年下半年供给缺口出来了，价格又涨回 28 块。

然后新一轮补栏开始。

\begin{chainflow}
\chainnode{猪价高}\chainarrow \chainnode{补栏}\chainarrow \chainnode{14 个月后集中出栏}\chainarrow \chainnode{价格暴跌}\chainarrow \chainnode{淘汰母猪}\chainarrow \chainnode{一年后供给不足}\chainarrow \chainnode{价格暴涨}
\end{chainflow}

\begin{egbox}{宋老板的两笔账}
2020 年：养殖完全成本约 15.6 元/公斤，售价 33 元/公斤。一头 120 公斤的猪，赚 \textbf{2088 元}。

2021 年 10 月：饲料涨价，成本变成 17.2 元/公斤，售价 10.8 元/公斤。同样一头 120 公斤的猪，亏 \textbf{768 元}。

同一个猪场，同一个人，同样的管理水平，一头猪的盈亏差了 2856 元。

中间他什么都没做错。他只是在 2020 年做了一个当时看完全正确、而且所有同行都在做的决定。

\end{egbox}

这套机制有个名字，叫\term{蛛网模型}{cobweb model}\index{088@蛛网模型}：生产决策依据的是上一期的价格，而产品真正出来的时候，价格已经被所有人的同一个决策改掉了。把价格和产量画在一张图上，轨迹会绕成一圈一圈的蛛网。

关键在于它收敛还是发散。如果供给对价格的反应比需求更敏感，蛛网就越绕越大，波动自我放大。生猪是典型的发散型。

往下推一步就明白了：\hlmark{时滞越长的品种，波动越猛。}鸡蛋的周期比猪短，大蒜和生姜比猪更短更疯。多晶硅料的产能建设周期是 18 到 24 个月，它的价格从 2022 年高点到 2024 年低点跌掉八成以上，用的是同一个模型。锂也是。

把这个模型放大到整个经济，就是库存周期。

\section{四个钟摆挂在一起}

\begin{tablewrap}
\begin{xltabular}{\linewidth}{>{\raggedright\arraybackslash}p{0.20\linewidth}rXX}
\toprule
\bthead{周期} & \bthead{长度} & \bthead{驱动力} & \bthead{盯什么} \\
\midrule
\textbf{库存（基钦）} & 3–4 年 & 订单波动沿产业链被层层放大 & PMI 新订单减产成品库存、工业企业产成品存货同比 \\
\textbf{设备资本开支（朱格拉）} & 7–11 年 & 设备寿命、技术换代、产能利用率 & 固定资产投资中的设备工器具购置、产能利用率 \\
\textbf{房地产建筑（库兹涅茨）} & 15–25 年 & 人口迁移、城镇化、家庭形成 & 新开工面积、土地成交、二手房挂牌量 \\
\textbf{债务长周期} & 50–75 年 & 债务对收入比的累积与清算 & 宏观杠杆率、利息支出占 GDP 比重 \\
\textbf{政治周期} & 看换届 & 选举与换届改变政策的松紧节奏 & 财政赤字率、专项债发行节奏、规划年份 \\
\bottomrule
\end{xltabular}
\end{tablewrap}

它们同时在走，互相干扰。所以你看到的从来不是一条干净的正弦曲线，是四五条不同频率的波叠加之后的那团毛线。

2008 年之所以那么惨，是因为库存周期、设备周期、地产周期和债务周期在同一个时间点全部转向下行。这种事情几十年才凑一次。

\section{周期为什么非存在不可}

\textbf{产能建设有时滞。}一座 12 英寸晶圆厂从动土到量产要 24 到 30 个月，一个铜矿从勘探到出矿七到十年。所有人在同一个价格信号下同时决定扩产，两年后同时投产。宋老板的猪场只是这件事最好懂的版本。

\textbf{信贷会加速一切。}抵押品升值，你能借到更多；借到更多，你买更多资产；资产更贵，抵押品又升值。这个正反馈叫\term{金融加速器}{financial accelerator}\index{028@金融加速器}。它在上行时把繁荣推得更高，在下行时同样对称地起作用，而且下行更快，因为催收比放贷容易。

\textbf{人会从众，而且是被制度逼的。}一个基金经理跑输两个季度要写报告解释，跑输四个季度可能要换人。"和大家一起错"的职业成本，远低于"一个人对"。凯恩斯那句话讲得最准：宁可循规蹈矩地失败，也不要标新立异地成功。

第三条最容易被当成道德问题。它不是。你把任何一个理性的人放进那个考核机制里，他都会做出同样的选择。

\bookbreak

\section{那根最长的钟摆}

短周期决定你今年好不好过，债务长周期决定你这一代人赶上的是什么剧本。

机制其实很朴素：加杠杆，资产涨，抵押品升值，于是借得更多。只要债务增速持续高于收入增速，这个循环就能自我强化很多年，看起来像永动机。

直到某个时点，利息负担吃掉了现金流。

\begin{egbox}{利息是怎么悄悄吃掉一家公司的}
一家公司年收入 3.2 亿，息税前利润 4100 万，负债 2.6 亿，平均融资成本 4.2\%。年利息 1092 万，利息保障倍数 3.8 倍。健康。

它每年借新钱扩张。三年后：负债 4.9 亿，收入 4.4 亿，息税前利润 5300 万。同期利率从 4.2\% 升到 6.5\%。

年利息变成 3185 万。利息保障倍数掉到 \textbf{1.66 倍}。

看它的报表：收入涨了 37\%，利润涨了 29\%，还在增长。但它离"利息吃掉全部经营利润"只剩不到一倍的空间。再来一次加息，或者一个季度的订单下滑，它就得开始卖资产。

这不是虚构的极端案例。过去三年很多中型企业走的就是这条路。

\end{egbox}

然后进入去杠杆。可用的手段只有四种：

\begin{tablewrap}
\begin{xltabular}{\linewidth}{>{\raggedright\arraybackslash}p{0.20\linewidth}XX}
\toprule
\bthead{手段} & \bthead{谁承担损失} & \bthead{后果} \\
\midrule
\textbf{紧缩（少花多还）} & 债务人 & 所有人一起省钱，收入同步下滑，债务对收入比可能不降反升 \\
\textbf{债务重组与违约} & 债权人 & 债权人的资产负债表受损，可能引发连锁 \\
\textbf{财富转移（税收与转移支付）} & 高财富群体 & 政治阻力最大，推进最慢 \\
\textbf{印钞与债务货币化} & 所有持币者 & 名义收入上升摊薄债务，代价是货币贬值和通胀 \\
\bottomrule
\end{xltabular}
\end{tablewrap}

四种手段的配比不同，结果完全不同。

全靠第一种，是 1930 年代的通缩螺旋：每个人都在还债，每个人的收入都是别人的支出，于是收入下降得比债务还快，杠杆率越去越高。全靠第四种，是货币贬值和恶性通胀。四者平衡、让名义增长率略高于名义利率、债务比重慢慢化掉而通胀不失控，达利欧管这叫\term{美丽的去杠杆}{beautiful deleveraging}\index{037@美丽的去杠杆}。

\begin{pullquote}{}
这四个旋钮怎么拧，从来不是经济学问题，是政治问题。它决定的是谁来承担损失，而这件事不由计算器决定。
\end{pullquote}

\section{中国站在哪一段}

不给预测，只摆位置。

房地产处在库兹涅茨长周期的调整段。这是十五到二十五年尺度的东西，不是一个季度的政策能翻转的，2026 年 8 月底出台的政策"组合拳"改变的是斜率和节奏，不是周期本身。

地方债务在化解，用的是展期、置换、降息，把短的换成长的、贵的换成便宜的。对照上面那张表，这是第一种和第二种手段的温和混合，避开了第四种。

人口结构过了拐点。这一条会同时压住地产周期和债务周期的分母。

产业升级和部分行业的产能过剩同时发生。光伏、锂电、一部分新能源车产业链，是朱格拉周期在同一时点被政策和资本一起推起来之后的自然结果。被叫做"内卷"的那个东西，用蛛网模型解释比用道德解释准确得多。

\begin{statrow}{3}
  \statitem{4.7}{\%}{中国 2026 年上半年 GDP 同比}
  \statitem{60}{万亿元}{2030 年社会消费品零售总额目标}
  \statitem{3.5}{\%}{5 年期以上 LPR（2026 年 8 月 20 日）}
\end{statrow}

商务部等 7 部门提出，到 2030 年社会消费品零售总额达到 60 万亿元左右。

这句话的信息量比数字本身大。\textbf{政策重心从投资转向消费，这本身就是周期位置的一个信号。}一个经济体处在投资周期的上行段，政策语言是"扩大有效投资"；处在投资周期的下行段和债务周期的去杠杆段，语言会变成"提振消费""扩大内需"。

官方文件的措辞变化，通常比统计数据更早。

\section{全球的账，和一个不能不提的尾部风险}

\begin{databox}{2026 数据}
国际金融协会（IIF）数据：截至 2026 年 3 月底，全球债务规模约 \textbf{353 万亿美元}，创历史新高，占全球 GDP 约 \textbf{305\%}。

IIF 点名的结构性推升力量有五个：人口老龄化、国防开支增加、能源安全与来源多元化、网络安全、AI 相关资本支出。

\end{databox}

这五样有一个共同点：都是不得不花的钱。

这一点让本轮债务扩张和过去几轮很不一样。过去的债务扩张主要由消费和地产驱动，那种债务在利率上行时是可以被压下去的：少买一套房，晚换一辆车，需求就萎缩了。

这一轮压不下去。养老金的支付义务不会因为国债收益率上行而减少；国防预算不会因为融资成本高就砍掉；算力如果不投，输掉的是一个国家未来二十年在产业链上的位置。

所以本轮债务对利率的敏感度比过去低。这句话听起来像好消息，其实是坏消息：\hlwarm{不敏感意味着即使利率很高，债务也停不下来。}

\begin{warnbox}{小心：2026 年的这个组合}
供给冲击（中东局势推高能源价格）＋ 高利率（美联储 3.50\%–3.75\%，日本央行 6 月加息至 1.0\%，欧洲央行 6 月三年来首次加息）＋ 高债务（353 万亿美元，占 GDP 305\%）。

债务可持续性说到底取决于一个不等式：名义利率 r 和名义增长率 g 谁大。r 小于 g，债务占 GDP 的比重会自己往下走；r 大于 g，就算一分钱新债都不借，比重也会自己往上爬。

滞胀是唯一一种能同时把这两边推向坏处的组合：通胀逼着央行抬高 r，停滞压低 g。

这不是预测。没有人知道下半年油价往哪走。这是在指出这个组合为什么危险。

\end{warnbox}

\begin{egbox}{r 和 g 差两个百分点是什么概念}
一个国家债务占 GDP 200\%，名义 GDP 增速 5\%，平均融资成本 3\%。就算财政收支完全平衡，债务比重每年也会自己下降约 3.8 个百分点：200\% × (3\% − 5\%) ÷ 1.05。

换个环境：增速掉到 2\%，融资成本升到 4.5\%。同样一分钱新债不借，比重每年自己上升约 4.9 个百分点。

两种情形里没有任何人多借一分钱，也没有任何政策失误。差别只来自两个数字谁大谁小。

这就是为什么各国政府对通胀的态度那么纠结：温和的通胀能抬高 g，帮着化债；失控的通胀会先把 r 顶上去，把好处全吐回来。

\end{egbox}

\begin{mythbox}{常见误解}
\mvfrow{误解}{"周期有固定长度，把历史数据一平均就能算出下一个拐点。"}
\mvfrow{事实}{那些长度是事后统计出来的平均值，方差大到没法用来做决策。美国战后最短的一轮扩张只有 12 个月，最长的一轮 128 个月（2009 到 2020）。拿平均值去推算下一次，等于拿"人均寿命 79 岁"去推算某一个具体的人哪天走。周期是存在的，节拍器是不存在的。}
\end{mythbox}

\begin{mythbox}{常见误解}
\mvfrow{误解}{"政府有那么多工具，早晚能把周期熨平。"}
\mvfrow{事实}{政策能改变周期的形状，削平最高点，垫高最低点，但改变不了它存在。更要紧的是，政策本身就是下一个周期的原因之一。2008 年之后的超低利率托住了资产价格，也养出了后来的高杠杆；2020 年的财政刺激避免了衰退，也贡献了 2021 到 2023 年的通胀。救助的代价通常有两种形式：一种是把问题挪到未来，一种是把问题挪给别人（储蓄者、后来买房的人、拿固定收入的人）。这两种代价都不会出现在救助当天的新闻标题里。}
\end{mythbox}

说清楚一件事：这一章没有任何内容可以用来择时，本书也不提供投资建议。判断位置和预测拐点是两回事，前者影响你该背多重的包，后者不影响任何事，因为它做不到。

\begin{takeaway}{本章一句话}
周期不是钟表，是四五个不同长度的钟摆挂在一起。你算不出下一个拐点，但你能判断自己站在哪一段：债务是在长还是在还，产能是在建还是在关，新进场的人是多还是少。

\begin{itemize}
  \item 库存 3–4 年，设备 7–11 年，地产 15–25 年，债务 50–75 年。同时在走，互相干扰。
  \item 时滞、信贷加速器、从众。三个机制里只有第三个看着像人性问题，其实是考核机制问题。
  \item 债务增速长期高于收入增速，终点是确定的，时间是不确定的。
  \item 去杠杆的四个旋钮怎么拧，决定谁来承担损失。这是政治问题。
\end{itemize}

\end{takeaway}

\begin{bridgebox}
\textbf{下一章}天气看完了。接下来只剩一个问题：这些东西跟你银行卡里那点钱，到底有什么关系？
\end{bridgebox}

